% CSE 715 — Computer Graphics
% Z-Buffer Algorithm (Chapter 10, Hidden Surfaces §10.1-10.2) — Complete Model Answers (2020–2024)
% University of Chittagong, 7th Semester B.Sc. Engineering
% Source: Schaum's Outline of Computer Graphics (Xiang & Plastock), Chapter 10
%
% Compile with: pdflatex ZBuffer_Solutions.tex (run twice for TOC)

\documentclass[12pt, a4paper]{article}

\usepackage[left=2.5cm, right=2.5cm, top=2.8cm, bottom=2.8cm, headheight=15pt]{geometry}
\usepackage{amsmath, amssymb}
\usepackage{tikz}
\usetikzlibrary{arrows.meta, positioning, shapes.geometric, calc, backgrounds, fit}
\usepackage{booktabs}
\usepackage{array}
\usepackage{xcolor}
\usepackage{enumitem}
\usepackage{fancyhdr}
\usepackage{titlesec}
\usepackage{mdframed}
\usepackage{tabularx}
\usepackage{hyperref}

\definecolor{sectionblue}{RGB}{10,60,110}
\definecolor{boxbg}{RGB}{242,247,255}
\definecolor{answerbg}{RGB}{240,255,240}
\definecolor{notesbg}{RGB}{255,250,230}

\mdfdefinestyle{solutionframe}{linecolor=sectionblue, backgroundcolor=boxbg, roundcorner=4pt, linewidth=1.4pt, innertopmargin=7pt, innerbottommargin=7pt, innerleftmargin=9pt, innerrightmargin=9pt}
\mdfdefinestyle{answerframe}{linecolor=green!55!black, backgroundcolor=answerbg, roundcorner=4pt, linewidth=1.4pt, innertopmargin=7pt, innerbottommargin=7pt, innerleftmargin=9pt, innerrightmargin=9pt}
\mdfdefinestyle{notesframe}{linecolor=orange!70!black, backgroundcolor=notesbg, roundcorner=4pt, linewidth=1pt, innertopmargin=5pt, innerbottommargin=5pt, innerleftmargin=8pt, innerrightmargin=8pt}
\newenvironment{solutionbox}{\begin{mdframed}[style=solutionframe]}{\end{mdframed}}
\newenvironment{answerbox}{\begin{mdframed}[style=answerframe]}{\end{mdframed}}
\newenvironment{notesbox}{\begin{mdframed}[style=notesframe]}{\end{mdframed}}

\pagestyle{fancy}
\fancyhf{}
\lhead{\small\color{sectionblue} CSE 715 --- Z-Buffer Algorithm (Ch10)}
\rhead{\small\color{sectionblue} Model Solutions (2020--2024)}
\cfoot{\small\thepage}
\renewcommand{\headrulewidth}{0.5pt}

\titleformat{\section}{\Large\bfseries\color{sectionblue}}{}{0em}{}[\vspace{-2pt}\color{sectionblue}\rule{\linewidth}{0.8pt}]
\titleformat{\subsection}{\large\bfseries\color{sectionblue!85!black}}{}{0em}{}

\begin{document}

\begin{center}
  {\LARGE\bfseries\color{sectionblue} CSE 715: Computer Graphics}\\[0.6em]
  {\Large\bfseries Z-Buffer Algorithm --- Chapter 10, Hidden Surfaces (\S10.1--10.2)}\\[0.3em]
  {\large Complete Model Answers (2020--2024)}\\[0.2em]
  {\normalsize University of Chittagong \quad|\quad 7\textsuperscript{th} Semester B.Sc.\ (Engg.) Examination}\\[0.3em]
  {\small Source: Schaum's Outline of Computer Graphics (Xiang \& Plastock), \S10.1--10.2 --- notation follows the book exactly ($Z_{\mathrm{buf}}(x,y)$, $Z_{\mathrm{back}}$)}
\end{center}

\vspace{0.3em}\noindent\rule{\linewidth}{1pt}\vspace{0.5em}
\tableofcontents
\newpage

%=====================================================================
\section{Theoretical Framework}
%=====================================================================
\begin{notesbox}
\textbf{Yield note:} re-verified against all 5 raw scanned papers 2026-07-06. The \textbf{full "max objects + 2$\times$2 pixel trace"} question is only \textbf{2021 (Q6d) and 2022 (Q6c)}. 2023 (Q7d) asks a lighter "differentiate Painter's vs.\ Z-buffer" instead, and 2024 (Q7d) asks "illustrate + advantages/disadvantages." \textbf{2020 has no Z-buffer question at all} --- that year's Ch10 content is only the separate "point-obscures-point via viewpoint" topic (Q8c), tracked independently.
\end{notesbox}

\subsection{Depth Comparison (\S10.1)}
\begin{answerbox}
Given two points $P_1(x_1,y_1,z_1)$ and $P_2(x_2,y_2,z_2)$ on the \emph{same projection line} (i.e.\ the same pixel), a depth comparison tells us which is in front. For an orthographic parallel projection onto the $xy$ plane with the viewing direction along $+z$ (viewer at $-z$, looking toward $+z$): \textbf{if $z_1<z_2$, then $P_1$ obscures $P_2$} (smaller $z$ = closer to the viewer). All hidden-surface algorithms in this chapter assume coordinates already in \textbf{normalized viewing coordinates}, with any perspective view first converted to an equivalent parallel one via the perspective-to-parallel transform $T_p$, so depth comparison always reduces to this simple $z_1$ vs.\ $z_2$ test.
\end{answerbox}

\subsection{The Z-Buffer Algorithm (\S10.2)}
\begin{answerbox}
An \textbf{image-space} method. Maintain a $Z$-buffer $Z_{\mathrm{buf}}(x,y)$ --- the depth of the closest point seen so far at each pixel $(x,y)$.
\begin{enumerate}[noitemsep]
  \item \textbf{Initialize:} set every pixel to the background color; set $Z_{\mathrm{buf}}(x,y) = Z_{\mathrm{back}}$ (the depth of the back clipping plane) for every pixel.
  \item \textbf{Scan-convert each (already-projected) polygon $P$ in the scene.} For every pixel $(x,y)$ inside $P$: compute $Z(x,y)$, the depth of $P$ at that pixel.
  \begin{itemize}[noitemsep]
    \item If $Z(x,y) < Z_{\mathrm{buf}}(x,y)$: set $Z_{\mathrm{buf}}(x,y)=Z(x,y)$ and set the pixel's color to $P$'s color at $(x,y)$ (this polygon is closer than anything seen so far at this pixel).
    \item Otherwise: leave the pixel and $Z_{\mathrm{buf}}$ unchanged (something already drawn is closer).
  \end{itemize}
\end{enumerate}
\textbf{Maximum number of objects representable:} \textbf{unlimited / arbitrarily large.} The algorithm requires $Z$-buffer memory proportional only to the number of \emph{pixels} on screen --- it processes polygons \emph{one at a time} and never needs to store all the scene's objects simultaneously, so the total object count in a scene is not bounded by the algorithm itself.
\textbf{Advantages:} simple to implement; handles an arbitrary/unlimited number of objects; no need to sort or resolve polygon-polygon depth ordering ambiguities (unlike Painter's algorithm); processes polygons independently, in any order.
\textbf{Disadvantages:} needs $Z$-buffer memory proportional to screen resolution (one depth value per pixel, in addition to the frame buffer); can suffer from depth-precision/aliasing artifacts; doesn't natively handle transparency (only the single nearest opaque surface per pixel is kept); no inherent anti-aliasing at silhouette edges.
\end{answerbox}

\subsection{Painter's Algorithm vs.\ Z-Buffer (for the differentiate-style question)}
\begin{answerbox}
\begin{center}
\begin{tabular}{p{4.3cm}p{5.3cm}p{5.3cm}}
\toprule
 & \textbf{Painter's Algorithm} & \textbf{Z-Buffer} \\
\midrule
Method & Depth-sort: paint polygons farthest-to-nearest by priority order & Per-pixel depth compare; keep closest seen so far \\
Ordering effort & Must sort by $z_{\max}$, then resolve overlap ambiguities (Tests 0--5); may need to split polygons & No sorting needed --- polygons processed in any order \\
Memory & No extra per-pixel depth storage, but needs polygon list \& ambiguity-resolution structures & Needs one $Z$-buffer entry per pixel \\
Ambiguity handling & Cyclic/overlapping polygons can force polygon splitting & None --- simple numeric compare always resolves it \\
Object count limit & Sorting cost grows with polygon count & Effectively unlimited (memory is per-pixel, not per-object) \\
\bottomrule
\end{tabular}
\end{center}
\end{answerbox}

%=====================================================================
\section{Examination 2022 --- Q6(c) \& 2021 --- Q6(d): Full Trace}
%=====================================================================
\begin{solutionbox}
\textbf{2022 Q6(c) [4.5 marks]:} What is the maximum number of objects that can be presented by using the Z-buffer algorithm? Using a 2$\times$2 pixel display, show how the Z-buffer algorithm would determine the color of each pixel for the given objects A and B in the figure (a box scene with a back clipping plane, objects A and B shown at different depths, pixel coordinates 0/1/2 on both axes).

\textbf{2021 Q6(d) [4.75 marks]:} Identical question structure with a different A/B depth arrangement in the figure.
\end{solutionbox}

\subsection*{Part 1 --- Maximum Number of Objects}
\begin{answerbox}
\textbf{Unlimited / arbitrarily large} --- see \S1.2. The $Z$-buffer's memory requirement depends only on screen resolution (one depth value per pixel), not on how many objects are in the scene, since objects are scan-converted and compared one at a time.
\end{answerbox}

\subsection*{Part 2 --- 2$\times$2 Pixel Trace Method}
\begin{notesbox}
The exact depth values are read off each year's own diagram (a 3D box showing objects A and B at different distances behind the display screen, with pixel coordinates $0,1,2$ on both axes of a $2\times2$ display). Apply the method below directly using your paper's specific A/B positions and depths --- the mechanics never change, only which of A/B is closer at each of the 4 pixels.
\end{notesbox}
\begin{solutionbox}
\textbf{Step 1 --- Initialize:} for all 4 pixels $(0,0),(1,0),(0,1),(1,1)$, set color = background, $Z_{\mathrm{buf}} = Z_{\mathrm{back}}$ (the depth of the rear clipping plane, the largest possible depth).

\textbf{Step 2 --- Scan-convert object A first (arbitrary order):} for every pixel that A's projection covers, read A's depth $Z_A(x,y)$ at that pixel. Since $Z_A < Z_{\mathrm{back}}$ everywhere A is drawn, update $Z_{\mathrm{buf}}(x,y)=Z_A(x,y)$ and set that pixel's color to A's color.

\textbf{Step 3 --- Scan-convert object B:} for every pixel B covers, read $Z_B(x,y)$ and compare against the current $Z_{\mathrm{buf}}(x,y)$:
\begin{itemize}[noitemsep]
  \item At any pixel covered by \emph{both} A and B: whichever of $Z_A,Z_B$ is \textbf{smaller} (nearer the viewer) wins --- that object's color is the final pixel color.
  \item At a pixel covered only by B (A didn't reach it): $Z_B < Z_{\mathrm{buf}}=Z_{\mathrm{back}}$, so B's color wins there.
  \item At a pixel covered by neither: the color stays background.
\end{itemize}

\textbf{Worked illustrative example} (representative numbers, method identical to either year's actual figure): suppose A occupies pixels $(0,0)$ and $(1,0)$ at depth $z=3$; B occupies pixels $(1,0)$ and $(1,1)$ at depth $z=2$; pixel $(0,1)$ is covered by neither.
\begin{center}
\begin{tabular}{ccccc}
\toprule
Pixel & Covered by & Depths seen & Winner ($\min z$) & Final color \\
\midrule
$(0,0)$ & A only & $z_A=3$ & A & color of A \\
$(1,0)$ & A and B & $z_A=3,\ z_B=2$ & B ($2<3$) & color of B \\
$(0,1)$ & neither & --- & --- & background \\
$(1,1)$ & B only & $z_B=2$ & B & color of B \\
\bottomrule
\end{tabular}
\end{center}
\textbf{Substitute your own paper's A/B pixel coverage and depth values into this same table structure} --- the answer is always "smallest $z$ wins at each pixel; background where neither object is drawn."
\end{solutionbox}

%=====================================================================
\section{Examination 2023 --- Q7(d): Painter's vs.\ Z-Buffer [2 marks]}
%=====================================================================
\begin{solutionbox}
\textbf{Question:} Differentiate between the Painter's algorithm and the Z-buffer algorithm.
\end{solutionbox}
\begin{answerbox}
See the comparison table in \S1.3. Core distinction to lead with: Painter's algorithm is a \textbf{depth-sort / priority} method that paints whole polygons in far-to-near order and must resolve sorting ambiguities (possibly splitting polygons); the Z-buffer is a \textbf{per-pixel numeric-comparison} method requiring no sorting at all, at the cost of needing a full-screen depth buffer in memory.
\end{answerbox}

%=====================================================================
\section{Examination 2024 --- Q7(d): Illustrate + Pros/Cons [2 marks]}
%=====================================================================
\begin{solutionbox}
\textbf{Question:} Illustrate the Z-buffer algorithm with necessary figures. What are its advantages and disadvantages?
\end{solutionbox}
\begin{answerbox}
\textbf{Illustration:} draw the initialize/scan-convert/compare-and-update flow from \S1.2 (a simple flowchart: Initialize $Z_{\mathrm{buf}}=Z_{\mathrm{back}}$ for all pixels $\to$ for each polygon, for each covered pixel, compute $Z(x,y)$ $\to$ if $Z(x,y)<Z_{\mathrm{buf}}(x,y)$ update buffer + pixel color $\to$ repeat for all polygons). \textbf{Advantages/disadvantages:} see \S1.2 --- lead with "handles unlimited objects, no sorting needed" as the main advantage and "needs full-screen depth memory, no native transparency/anti-aliasing" as the main disadvantage.
\end{answerbox}

%=====================================================================
\section{Quick Summary}
%=====================================================================
\begin{center}
\renewcommand{\arraystretch}{1.3}
\begin{tabular}{p{4cm}|p{2.2cm}|p{7.3cm}}
\toprule
\textbf{Pattern} & \textbf{Years} & \textbf{Method} \\
\midrule
Max objects + 2$\times$2 pixel trace (full numeric) & 2021, 2022 & Unlimited objects (memory is per-pixel); trace: smallest $z$ wins at each pixel, background where nothing drawn \\
Differentiate Painter's vs.\ Z-buffer & 2023 & Depth-sort+ambiguity-resolution vs.\ per-pixel numeric compare \\
Illustrate + advantages/disadvantages & 2024 & Flowchart of init/compare/update; unlimited objects vs.\ full-screen memory cost \\
\bottomrule
\end{tabular}
\end{center}
\begin{notesbox}
\textbf{Exam-day takeaway:} whatever form it takes, the answer always centers on the same fact --- \textbf{Z-buffer compares depths per pixel, keeps the smallest, and needs one depth value per pixel rather than a bound on object count.} Memorize that one sentence and you can answer any of the four question shapes above.
\end{notesbox}

\end{document}
